\section{Geometry of the layers: nesting, chains and polar rules}\label{lg:sec}

This section proves the geometric facts about the convex layers of a rotation-invariant finite point set that the
layered encoders rely on, together with the elementary identities behind them. In the statements about the layers
(Lemmas~\ref{lg:lem:layer} and~\ref{lg:lem:halfplane}, Theorems~\ref{lg:thm:nesting}, \ref{lg:thm:radius}
and~\ref{lg:thm:chain}) and in Lemma~\ref{lg:lem:centre}, $S$ is a finite set of points that is mapped into itself by
the counterclockwise rotation $\rho$ by $2\pi/m$ ($m\ge1$) about a point $c$; in all of them except
Lemma~\ref{lg:lem:centre}, also $c\notin S$ and $S$ has no $K$ points in convex position for some $K\le 2m$.
Theorem~\ref{lg:thm:fan}, the only other statement with a set $S$, involves no centre and no rotation. The section
contains:
\begin{itemize}
\item \emph{rotations and regular polygons} (\S\ref{lg:ssec:rot}): iterates of a rotation, the orientation of three
  vertices of a regular polygon (Lemmas~\ref{lg:lem:orient-rot}, \ref{lg:lem:trig}, \ref{lg:lem:regular},
  \ref{lg:lem:next-prev}) and the position of the vertices relative to the tangent at a vertex
  (Lemmas~\ref{lg:lem:lin-rot}, \ref{lg:lem:lin-orbit});
\item \emph{convexity} of the sets $\{q:0\le\orient(a,b,q)\}$ (closed half-planes when $a\ne b$) and of closed disks
  (\S\ref{lg:ssec:convex});
\item \emph{nesting, centre and radii} (\S\ref{lg:ssec:nest}): if $m\ge2$ and $S$ is in general position, every point
  of a later layer lies strictly to the left of every edge $x\to\rho(x)$ of an earlier layer
  (Theorem~\ref{lg:thm:nesting}); if $m\ge2$ is even and $S$ is in general position and has a point other than $c$,
  then $c\notin S$ (Lemma~\ref{lg:lem:centre}); the points of a layer and of all later layers are no farther from $c$
  than the vertices of that layer (Theorem~\ref{lg:thm:radius});
\item \emph{fans and chains} (\S\ref{lg:ssec:chain}): a pivot followed by $n\ge3$ chain points with positive fan and
  turn orientations, all in an open half-plane whose boundary line passes through the pivot, gives $n+1$ points in
  convex position (Theorem~\ref{lg:thm:fan}; by its proof, the pivot and the chain points themselves); if $m\ge3$ and
  $S$ is in general position, all other points of the layers $\ge i$ lie strictly on the inner side of the tangent at
  a vertex of layer $i$ (Lemma~\ref{lg:lem:halfplane}); hence, if moreover $K\ge4$, no vertex of layer $i$ is the pivot
  of a sequence of $K-1$ points of the layers $\ge i$, all different from it, with positive fan and turn orientations
  (Theorem~\ref{lg:thm:chain}), which is the geometric content of the chain clauses (see the paragraph
  \emph{Encoders} below);
\item \emph{polar rules} (\S\ref{lg:ssec:polar}): the orientation of three points given in polar coordinates about $c$
  (Lemma~\ref{lg:lem:polar}) and the angular rules A and B (Lemmas~\ref{lg:lem:ruleA}, \ref{lg:lem:ruleB}).
\end{itemize}

\paragraph{Sources.} Every Lean declaration that heads a statement of this section is in the file
\dirn{LAYERLEAN/LayerLemma.lean} (namespace \lean{ErdosSzekeres.LayerLemma}, 1148 lines, SHA-256
\texttt{afb8542dceb0}\dots), lines 465--1032; its README is \dirn{LAYERLEAN/README.md} (SHA-256
\texttt{9eaec241a1d2}\dots). The declarations of the same file that \S\ref{lg:ssec:notation} recalls lie outside
these lines and are restated in the section on the layer lemma. The Lean proofs of the declarations stated here use the following
declarations of the main library; the foundations section restates them and proves the lemmas among them.
\begin{itemize}
\item The definitions \lean{Point}, \lean{GeneralPosition}, \lean{ConvexPosition} and \lean{HasConvexNGon}
  of\newline \dirn{ErdosSzekeres/Definitions.lean} (SHA-256 \texttt{2425500eefd1}\dots).
\item The orientation \lean{ErdosSzekeres.orient} and the lemma \lean{ErdosSzekeres.orient_cycle},\newline
  in \dirn{ErdosSzekeres/Orientation.lean}, lines 9 and 15 (SHA-256 \texttt{f67955ac7091}\dots).
\item The lemma \lean{ErdosSzekeres.generalPosition_iff_orient_ne_zero},\newline
  in \dirn{ErdosSzekeres/OrientationGeometry.lean:66} (SHA-256 \texttt{c7db9bbbb65e}\dots).
\item The lemma \lean{ErdosSzekeres.Suk.convexPosition_of_positive_triples},\newline
  in \dirn{ErdosSzekeres/Suk/OrientationClique.lean:9} (SHA-256 \texttt{cf63d5219b47}\dots).
\end{itemize}
The written proofs use the results listed under \emph{Dependencies} below, among them Lemma~\ref{fd:lem:mono}(c)
(Lean \lean{HasConvexNGon.mono}, \dirn{ErdosSzekeres/Basic.lean:20}), which the Lean proofs do not use.

The file was checked with Lean~4.24.0 and Mathlib \texttt{f897ebcf} by \texttt{lake env lean} from the repository root;
it is not part of the default \texttt{lake build} target. The log \dirn{LAYERLEAN/logs/build.log} (exit code~0, no
messages) names no file hash; it was committed with the first version of the file (454 lines, commit
\texttt{6429e9a70}) and was not updated by the five later commits that changed the file, so it is not a record of the
present file, and we do not use it. The records for the present file are the review report
\dirn{LEANEXT/review/REPORT.md} of 4 November 2026, which reports exit~0 and no messages for the reviewed file, whose
SHA-256 it gives as \texttt{afb8542d}\dots{} (1148 lines), and the axiom re-check of 6 November 2026 (whose output is
named below), which compiled
\dirn{LAYERLEAN/logs/axioms_recheck_20261106/LayerLemmaAxioms.lean}, whose first 1148 lines are this file, with exit
status~0; the files \dirn{README.txt} and \dirn{run.sh} of that directory record the toolchain and the command given
above. The file contains no \lean{sorry}, \lean{admit}, \lean{axiom} or \lean{native_decide}.

\paragraph{Axioms.} The log \dirn{LAYERLEAN/logs/axioms.log} (SHA-256 \texttt{3c4b7aac502e}\dots) prints the axioms of
16 of the 23 theorems of this section: all except \lean{rotAbout_iterate}, \lean{orient_rotAbout},
\lean{convex_orient_nonneg}, \lean{orient_eq_lin}, \lean{lin_rotAbout}, \lean{lin_orbit_pos} and
\lean{orient_next_prev_pos}. The review log \dirn{LEANEXT/review/axioms.log} of 4 November 2026 (SHA-256
\texttt{386c4f578725}\dots) prints all 23. It covers 59 declarations of the reviewed file (SHA-256
\texttt{afb8542d}\dots, the file above), namely all its \texttt{theorem} declarations except the 7 marked
\texttt{@[simp]}; the file has 66 in all, and \dirn{LEANEXT/review/REPORT.md} speaks of ``all 59 theorem
declarations''. In each of the two logs, every theorem of this section that it prints depends only on the standard
axioms \lean{propext}, \lean{Classical.choice} and \lean{Quot.sound}. The re-check
\dirn{LAYERLEAN/logs/axioms_recheck_20261106/out.txt} (SHA-256 \texttt{a59c3816cbe7}\dots) gives the same result for
the 11 of them that it lists.

\paragraph{Rigor.} Each lemma or theorem below whose head carries a Lean name is \emph{Lean-checked}, with the axioms
just listed; none of them takes an unproved proposition as a hypothesis, and none is a finite computation. Each definition
whose head carries a Lean name restates that Lean definition. The proofs written
here are \emph{written proofs (not refereed)}. They follow the structure of the Lean proofs. Where Lean checks an
algebraic identity or inequality with an automatic tactic (such as \lean{ring}, \lean{linear_combination},
\lean{nlinarith} or \lean{positivity}), the written proof derives it by hand; every other difference from the Lean
proof is stated in a remark or in a parenthetical note in the proof. Statements tagged ``written step'' are not Lean
declarations, and they are proved here. Lemma~\ref{lg:lem:layer} collects the opening steps of the Lean proofs of
Theorems~\ref{lg:thm:nesting} and~\ref{lg:thm:radius} and Lemma~\ref{lg:lem:halfplane}. Lemma~\ref{lg:lem:rotfacts} collects
elementary facts about rotations that the Lean proofs obtain differently (through the complex form of the rotation,
by direct computation, or through minimal periods). Lemma~\ref{lg:lem:cramer} is a general identity of which Lean
checks only the three instances used, each by \lean{ring}. Corollary~\ref{lg:cor:period} restates
Lemma~\ref{lg:lem:rot-iter} for $\rho$ and adds $\rho^m=\mathrm{id}$, which replaces a minimal-period argument of the
Lean proof of Lemma~\ref{lg:lem:halfplane}.

\paragraph{Encoders.} Several results below are the geometric facts behind parts of the project's layered encoders:
Theorem~\ref{lg:thm:nesting} behind the nesting clauses, Theorem~\ref{lg:thm:chain} behind the chain clauses, and
Lemmas~\ref{lg:lem:ruleA} and~\ref{lg:lem:ruleB}, with the weak radius order of Theorem~\ref{lg:thm:radius} that
rule~B needs, behind the angular rules A and B, which the encoders apply to triples of points. Only the geometric
facts are stated and proved here. That the
encoders' programs emit exactly these clause families, and which rule they apply to which triple, are not part of
\dirn{LAYERLEAN/LayerLemma.lean}; \dirn{LAYERLEAN/README.md} lists both under ``What is not in Lean'' and gives their
status as ``written proof, review and tests''. For the encoders of the $k=8$ bounds the code-to-clause step is
recorded as gap~(G1) of \dirn{WRITEUP/note/k8_soundness.tex}: it rests on a written proof (not refereed) by reading
the code, with the Python interpreter trusted to run the programs as written. For encoder~B at $m=6$ that note also
records a second route on which the step is Lean-checked for a Lean definition of the formula (\dirn{ENCBLEAN}). On
that route the note lists, under gap~(G2), the trust items (E1)--(E6), which are not Lean-checked: among them, that the
texts printed from this definition are the $64$ cube files that carry the cake\_lpr verdicts is (E1), evidence from
running compiled code and not a Lean theorem, and those $64$ verdicts and their provenance are covered by (E2).

\paragraph{Dependencies.} The section uses the conventions and facts of the foundations section
(Definitions~\ref{fd:def:point}, \ref{fd:def:orient}, \ref{fd:def:gp}, \ref{fd:def:cp}, \ref{fd:def:ngon},
\ref{fd:def:inv}, \ref{fd:def:rot} and Lemmas~\ref{fd:lem:orient}, \ref{fd:lem:collinear}, \ref{fd:lem:mono},
\ref{fd:lem:pos-triples}) and, from the section on the layer lemma, Definitions~\ref{ll:def:hv} and~\ref{ll:def:peel},
Lemmas~\ref{ll:lem:hv}, \ref{ll:lem:antitone}, \ref{ll:lem:inv-peel} and Theorem~\ref{ll:thm:core-rot}. These,
except Lemma~\ref{fd:lem:pos-triples}, which is used only in Step~3 of the proof of Theorem~\ref{lg:thm:fan}, are
recalled in \S\ref{lg:ssec:notation}.

\subsection{Notation and the results used}\label{lg:ssec:notation}

\paragraph{Vectors.} Points are elements of $\R^2$, written $p=(p_0,p_1)$ (Definition~\ref{fd:def:point}; Lean
\lean{Point}). For $X,Y\in\R^2$ put
\[
  \det(X,Y)=X_0Y_1-X_1Y_0,\qquad \langle X,Y\rangle=X_0Y_0+X_1Y_1,\qquad \abs{X}^2=\langle X,X\rangle,\qquad
  JX=(-X_1,X_0).
\]
So $J$ is the counterclockwise quarter turn, $\det$ is bilinear and antisymmetric, and $\det(X,Y)=\langle JX,Y\rangle$.
In particular
\begin{equation}\label{lg:eq:J}
\begin{gathered}
  \det(X,X)=0,\qquad \det(X,JX)=\abs{X}^2,\qquad \det(JX,X)=-\abs{X}^2,\\
  \det(JX,JX)=0,\qquad \langle X,JX\rangle=0 .
\end{gathered}
\end{equation}

\paragraph{Orientation.} We use the following facts about $\orient$ (Definition~\ref{fd:def:orient}).
\begin{itemize}
\item[(O1)] $\orient(a,b,d)=(b_0-a_0)(d_1-a_1)-(b_1-a_1)(d_0-a_0)=\det(b-a,\,d-a)$.
\item[(O2)] $\orient$ is invariant under cyclic permutations of its three arguments (Lemma~\ref{fd:lem:orient}(b)).
  Hence $\orient(a,b,o)=\orient(o,a,b)=\det(a-o,\,b-o)$ for all points $a,b,o$.
\item[(O3)] $\orient(a,b,d)=0$ whenever two of $a,b,d$ coincide. Indeed, by (O1), $\orient(a,a,d)=\det(0,d-a)=0$,
  $\orient(a,b,a)=\det(b-a,0)=0$ and $\orient(a,b,b)=\det(b-a,b-a)=0$.
\item[(O4)] For fixed $a,b$ the map $q\mapsto\orient(a,b,q)=\det(b-a,\,q-a)$ is affine: if $u+v=1$ then
  $\orient(a,b,uq_1+vq_2)=u\orient(a,b,q_1)+v\orient(a,b,q_2)$, because $(uq_1+vq_2)-a=u(q_1-a)+v(q_2-a)$ and $\det$ is
  linear in its second argument.
\end{itemize}

\paragraph{Rotations.} For $\theta\in\R$ let $R_\theta$ be the matrix with rows $(\cos\theta,-\sin\theta)$ and
$(\sin\theta,\cos\theta)$, so that $R_\theta X=\cos\theta\,X+\sin\theta\,JX$. Definition~\ref{fd:def:rot} reads
\[
  \operatorname{rotAbout}(c,\theta)(p)=c+R_\theta(p-c)
  =\bigl(c_0+\cos\theta\,(p_0-c_0)-\sin\theta\,(p_1-c_1),\;c_1+\sin\theta\,(p_0-c_0)+\cos\theta\,(p_1-c_1)\bigr).
\]
For a map $f$ and $j\in\N$, $f^{j}$ is the $j$-fold iterate: $f^{0}=\mathrm{id}$ and $f^{j+1}=f\circ f^{j}$
($=f^{j}\circ f$; Lean \texttt{f\^{}[j]}). For $m\in\N$ the angle $2\pi/m$ is the real number $2\pi/(m:\R)$ (Lean
\texttt{2 * Real.pi / m}, with $m$ cast to $\R$; Lean's convention would give $0$ for $m=0$, but every statement of
this section in which $2\pi/m$ occurs assumes $m\ge1$).
Whenever $c$ and $m$ are given we write
\[
  \rho=\operatorname{rotAbout}(c,2\pi/m).
\]
A natural-number subtraction such as $m-1$ or $K-1$ is truncated at $0$ in Lean; in every use below the minuend is at
least the subtrahend, so it is the ordinary difference.

\paragraph{Sine and cosine.} Besides the values at $0$, $\pi$ and $2\pi$, the addition, subtraction and double-angle
formulas, $\sin^2+\cos^2=1$, the oddness of the sine and the $2\pi$-periodicity, we use two standard facts about the real sine
and cosine (see, e.g., W.~Rudin, \emph{Principles of Mathematical Analysis}, 3rd ed., Chapter~8; in Mathlib
\lean{Real.sin_pos_of_pos_of_lt_pi}, \lean{Real.cos_le_one} and \lean{Real.cos_eq_one_iff}):
\begin{itemize}
\item[(S1)] $\sin\theta>0$ for $0<\theta<\pi$; hence $\sin\theta\ge0$ for $0\le\theta\le\pi$, as $\sin0=\sin\pi=0$;
\item[(S2)] $\cos\theta\le1$ for every $\theta\in\R$, and $\cos\theta=1$ if and only if $\theta\in2\pi\Z$.
\end{itemize}

\paragraph{Hypotheses on $S$.} All point sets are finite (Lean \texttt{Finset Point}). ``$\rho$ maps $S$ into $S$'' is the
Lean hypothesis \texttt{Invariant (rotAbout c (2 * Real.pi / m)) S}, i.e.\ $\rho(p)\in S$ for all $p\in S$
(Definition~\ref{fd:def:inv}); by induction on $i$ it implies $\rho^{i}(p)\in S$ for all $p\in S$ and $i\in\N$.
``$S$ is $K$-free'' means $\neg$\,\texttt{HasConvexNGon K S}: no $K$-element subset of $S$ is in convex position
(Definitions~\ref{fd:def:cp}, \ref{fd:def:ngon}). By Lemma~\ref{fd:lem:mono}(c), every subset of a $K$-free set is
$K$-free. ``General position'' is Definition~\ref{fd:def:gp}; by Lemma~\ref{fd:lem:collinear}, $S$ is in general
position if and only if $\orient(a,b,d)\ne0$ for all pairwise distinct $a,b,d\in S$.

\paragraph{Convexity.} A set $C\subseteq\R^2$ is convex (Mathlib \lean{Convex}) if $uy_1+vy_2\in C$ whenever
$y_1,y_2\in C$, $u,v\ge0$ and $u+v=1$. The convex hull $\conv(A)$ is the smallest convex set containing $A$; so
$A\subseteq C$ with $C$ convex implies $\conv(A)\subseteq C$ (Mathlib \lean{convexHull_min}).

\paragraph{Layers.} $\hv(X)=\{x\in X: x\notin\conv(X\setminus\{x\})\}\subseteq X$ (Definition~\ref{ll:def:hv}), and
$\peel(S,0)=S$, $\peel(S,n+1)=\peel(S,n)\setminus\hv(\peel(S,n))$ (Definition~\ref{ll:def:peel}). For a map $f$, a
number $k\in\N$ and a point $x$, Lean's \texttt{orbitFS f k x} (\dirn{LAYERLEAN/LayerLemma.lean:48}) is the finite set
$\{f^{j}(x):0\le j<k\}$. We use the following results of the section on the layer lemma.
\begin{itemize}
\item[(L1)] For every finite $X$: $X\subseteq\conv(\hv(X))$ (Lemma~\ref{ll:lem:hv}(c), \lean{subset_convexHull_hullVerts}).
\item[(L2)] If $n\le n'$ then $\peel(S,n')\subseteq\peel(S,n)$ (Lemma~\ref{ll:lem:antitone}); in particular
  $\peel(S,n)\subseteq\peel(S,0)=S$.
\item[(L3)] If $0<m$, $c\notin S$ and $\rho$ maps $S$ into $S$, then $\rho$ maps $\peel(S,n)$ into $\peel(S,n)$ for every
  $n$ (Lemma~\ref{ll:lem:inv-peel}, \lean{invariant_peel}).
\item[(L4)] If $0<m$, $K\le 2m$, $c\notin S$, $\rho$ maps $S$ into $S$, $S$ is $K$-free and $x\in\hv(S)$, then
  $\hv(S)=\{\rho^{j}(x):0\le j<m\}$ (Theorem~\ref{ll:thm:core-rot}, \lean{hullVerts_eq_orbitFS_rotAbout}).
\end{itemize}

The Lean proofs of Theorems~\ref{lg:thm:nesting} and~\ref{lg:thm:radius} and of Lemma~\ref{lg:lem:halfplane} each
begin with the same few steps; we collect them once.

\begin{lemma}[layers of a layered configuration; written step]\label{lg:lem:layer}
Let $m,K\in\N$ with $0<m$ and $K\le 2m$, let $S$ be a finite set and $c$ a point with $c\notin S$, let
$\rho=\operatorname{rotAbout}(c,2\pi/m)$ map $S$ into $S$, and let $S$ be $K$-free. Let $n\in\N$, put
$L=\hv(\peel(S,n))$, and let $x\in L$. Then:
\begin{enumerate}[label=(\alph*)]
\item $L=\{\rho^{j}(x):0\le j<m\}$; in particular $\rho^{j}(x)\in L$ for $0\le j<m$;
\item $L\subseteq\peel(S,n)\subseteq S$, and $x\neq c$;
\item $\peel(S,n')\subseteq\conv(L)$ for every $n'\ge n$;
\item $\peel(S,n')\cap L=\emptyset$ for every $n'>n$.
\end{enumerate}
\end{lemma}

\begin{proof}
By (L2), $\peel(S,n)\subseteq S$. Hence $c\notin\peel(S,n)$, and $\peel(S,n)$ is $K$-free (a $K$-element subset of
$\peel(S,n)$ in convex position would be one of $S$). By (L3), $\rho$ maps $\peel(S,n)$ into itself. So (L4) applies to
the finite set $\peel(S,n)$ and its hull vertex $x$, and gives (a).
(b) $L\subseteq\peel(S,n)$ by the definition of $\hv$, and $\peel(S,n)\subseteq S$; since $x\in L\subseteq S$ and
$c\notin S$, $x\ne c$.
(c) If $n'\ge n$ then $\peel(S,n')\subseteq\peel(S,n)\subseteq\conv(\hv(\peel(S,n)))=\conv(L)$ by (L2) and (L1).
(d) If $n'>n$ then $n+1\le n'$, so $\peel(S,n')\subseteq\peel(S,n+1)=\peel(S,n)\setminus L$ by (L2) and the definition
of $\peel$.
\end{proof}

\subsection{Rotations and regular polygons}\label{lg:ssec:rot}

\begin{lemma}[rotation facts; written step]\label{lg:lem:rotfacts}
For all points $c,p$ and all $\alpha,\beta,\theta\in\R$:
\begin{enumerate}[label=(\alph*)]
\item $\operatorname{rotAbout}(c,0)(p)=p$;
\item $\operatorname{rotAbout}(c,\alpha)\bigl(\operatorname{rotAbout}(c,\beta)(p)\bigr)=\operatorname{rotAbout}(c,\alpha+\beta)(p)$;
\item $\operatorname{rotAbout}(c,2\pi)(p)=p$;
\item $\abs{\operatorname{rotAbout}(c,\theta)(p)-c}^2=\abs{p-c}^2$.
\end{enumerate}
\end{lemma}

\begin{proof}
Put $u=p-c$, so that $\operatorname{rotAbout}(c,\theta)(p)=c+R_\theta u$.
(a), (c): $\cos0=\cos2\pi=1$ and $\sin0=\sin2\pi=0$, so $R_0u=R_{2\pi}u=u$.
(b) $\operatorname{rotAbout}(c,\beta)(p)-c=R_\beta u$, hence
$\operatorname{rotAbout}(c,\alpha)(\operatorname{rotAbout}(c,\beta)(p))=c+R_\alpha R_\beta u$. The product
$R_\alpha R_\beta$ has rows $(\cos\alpha\cos\beta-\sin\alpha\sin\beta,\;-(\sin\alpha\cos\beta+\cos\alpha\sin\beta))$ and
$(\sin\alpha\cos\beta+\cos\alpha\sin\beta,\;\cos\alpha\cos\beta-\sin\alpha\sin\beta)$, which is $R_{\alpha+\beta}$ by
the addition formulas $\cos(\alpha+\beta)=\cos\alpha\cos\beta-\sin\alpha\sin\beta$ and
$\sin(\alpha+\beta)=\sin\alpha\cos\beta+\cos\alpha\sin\beta$.
(d) $\abs{R_\theta u}^2=(\cos\theta\,u_0-\sin\theta\,u_1)^2+(\sin\theta\,u_0+\cos\theta\,u_1)^2
=(\cos^2\theta+\sin^2\theta)(u_0^2+u_1^2)=\abs{u}^2$; the mixed terms cancel.
\end{proof}

\begin{lemma}[iterates of the rotation about a point; \leanref{ErdosSzekeres.LayerLemma.rotAbout_iterate}{LAYERLEAN/LayerLemma.lean:465}]\label{lg:lem:rot-iter}
For every point $c$, every $\alpha\in\R$ and every $j\in\N$,
\[
  \bigl(\operatorname{rotAbout}(c,\alpha)\bigr)^{j}=\operatorname{rotAbout}(c,\,j\alpha)
\]
as maps $\R^2\to\R^2$, where $j$ is cast to $\R$.
\end{lemma}

\begin{proof}
Induction on $j$. For $j=0$ the left side is the identity by definition of the iterate, and the right side is
$\operatorname{rotAbout}(c,0)=\mathrm{id}$ by Lemma~\ref{lg:lem:rotfacts}(a). If the claim holds for $j$, then
\[
  \bigl(\operatorname{rotAbout}(c,\alpha)\bigr)^{j+1}=\operatorname{rotAbout}(c,\alpha)\circ\operatorname{rotAbout}(c,j\alpha)
  =\operatorname{rotAbout}(c,\alpha+j\alpha)=\operatorname{rotAbout}(c,(j+1)\alpha)
\]
by Lemma~\ref{lg:lem:rotfacts}(b). (Lean argues instead with the complex form: in the complex coordinate $z$, with
$\gamma=c_0+ic_1$, the map $\operatorname{rotAbout}(c,\alpha)$ is $z\mapsto\gamma+e^{i\alpha}(z-\gamma)$
(\lean{rotAbout_eq_crot}, line 302), whose $j$-th iterate is $z\mapsto\gamma+(e^{i\alpha})^{j}(z-\gamma)$
(\lean{crot_iterate}, line 267), and $(e^{i\alpha})^{j}=e^{ij\alpha}$.)
\end{proof}

\begin{corollary}[$\rho^m$ is the identity; written step]\label{lg:cor:period}
Let $c$ be a point, $m\in\N$ with $0<m$, and $\rho=\operatorname{rotAbout}(c,2\pi/m)$. Then $\rho^{j}=
\operatorname{rotAbout}(c,2\pi j/m)$ for every $j\in\N$, and $\rho^{m}(p)=p$ for every point $p$.
\end{corollary}

\begin{proof}
By Lemma~\ref{lg:lem:rot-iter}, $\rho^{j}=\operatorname{rotAbout}(c,j\cdot(2\pi/m))$. For $j=m$ the angle is
$m\cdot(2\pi/m)=2\pi$ because $(m:\R)\ne0$, and $\operatorname{rotAbout}(c,2\pi)=\mathrm{id}$ by
Lemma~\ref{lg:lem:rotfacts}(c).
\end{proof}

\begin{lemma}[orientation of a regular-polygon triple; \leanref{ErdosSzekeres.LayerLemma.orient_rotAbout}{LAYERLEAN/LayerLemma.lean:472}]\label{lg:lem:orient-rot}
For all points $c,x$ and all $\alpha,\beta\in\R$,
\[
  \orient\bigl(x,\operatorname{rotAbout}(c,\alpha)(x),\operatorname{rotAbout}(c,\beta)(x)\bigr)
  =\bigl((x_0-c_0)^2+(x_1-c_1)^2\bigr)\bigl((\cos\alpha-1)\sin\beta-\sin\alpha\,(\cos\beta-1)\bigr).
\]
\end{lemma}

\begin{proof}
Put $u=x-c$, $b=\operatorname{rotAbout}(c,\alpha)(x)$ and $d=\operatorname{rotAbout}(c,\beta)(x)$. Then
$b-x=R_\alpha u-u=(\cos\alpha-1)u+\sin\alpha\,Ju$ and $d-x=(\cos\beta-1)u+\sin\beta\,Ju$. By (O1), bilinearity of
$\det$ and~\eqref{lg:eq:J},
\begin{align*}
  \orient(x,b,d)=\det(b-x,d-x)
  &=(\cos\alpha-1)(\cos\beta-1)\det(u,u)+(\cos\alpha-1)\sin\beta\,\det(u,Ju)\\
  &\quad+\sin\alpha\,(\cos\beta-1)\det(Ju,u)+\sin\alpha\sin\beta\,\det(Ju,Ju)\\
  &=\bigl((\cos\alpha-1)\sin\beta-\sin\alpha\,(\cos\beta-1)\bigr)\abs{u}^2,
\end{align*}
and $\abs{u}^2=(x_0-c_0)^2+(x_1-c_1)^2$. No relation between $\cos$ and $\sin$ is used: this is a polynomial identity in
$\cos\alpha,\sin\alpha,\cos\beta,\sin\beta$ and the coordinates, which is how Lean proves it (\lean{ring}).
\end{proof}

\begin{lemma}[a product formula for the orientation factor; \leanref{ErdosSzekeres.LayerLemma.trig_layer}{LAYERLEAN/LayerLemma.lean:480}]\label{lg:lem:trig}
For all $s,t\in\R$,
\[
  (\cos 2t-1)\sin 2s-\sin 2t\,(\cos 2s-1)=4\sin t\,\sin s\,\sin(s-t).
\]
\end{lemma}

\begin{proof}
By the double-angle formulas, $\cos2t-1=-2\sin^2t$, $\cos2s-1=-2\sin^2s$, $\sin2s=2\sin s\cos s$ and
$\sin2t=2\sin t\cos t$. Hence the left side equals
\[
  -4\sin^2t\,\sin s\cos s+4\sin t\cos t\,\sin^2s=4\sin t\,\sin s\,(\sin s\cos t-\cos s\sin t)=4\sin t\,\sin s\,\sin(s-t).
\]
(Lean rewrites with $\cos2x=2\cos^2x-1$ and $\sin2x=2\sin x\cos x$ for $x=s,t$ and with the subtraction formula for
$\sin(s-t)$, and closes the remaining polynomial identity by \lean{linear_combination} with $\sin^2+\cos^2=1$ for $s$
and for $t$.)
\end{proof}

\begin{lemma}[vertices lie weakly left of each edge; \leanref{ErdosSzekeres.LayerLemma.orient_regular_nonneg}{LAYERLEAN/LayerLemma.lean:488}]\label{lg:lem:regular}
Let $c,x$ be points, $m\in\N$ with $0<m$, $j\in\N$ with $j\le m$, and $\rho=\operatorname{rotAbout}(c,2\pi/m)$. Then
\[
  0\le\orient\bigl(x,\rho(x),\rho^{j}(x)\bigr).
\]
\end{lemma}

\begin{proof}
By Lemma~\ref{lg:lem:rot-iter}, $\rho^{j}(x)=\operatorname{rotAbout}(c,\beta)(x)$ with $\beta=j\cdot(2\pi/m)$. By
Lemma~\ref{lg:lem:orient-rot} with $\alpha=2\pi/m$,
\[
  \orient(x,\rho(x),\rho^{j}(x))=\abs{x-c}^2\,F,\qquad F=(\cos\alpha-1)\sin\beta-\sin\alpha\,(\cos\beta-1).
\]
The first factor is $\ge0$; we show $F\ge0$. If $j=0$ then $\beta=0$, so $\sin\beta=0=\cos\beta-1$ and $F=0$. Let
$1\le j\le m$ and put $t=\pi/m$. Then $\alpha=2t$ and $\beta=2(jt)$, so Lemma~\ref{lg:lem:trig} with $s=jt$ gives
\[
  F=4\sin t\,\sin(jt)\,\sin(jt-t).
\]
Since $m\ge1$, $0<t\le\pi$; since $j\le m$, $0\le jt=j\pi/m\le\pi$; and since $j\ge1$, $0\le jt-t=(j-1)t\le jt\le\pi$.
By (S1) the sine is $\ge0$ on $[0,\pi]$, so all three factors are $\ge0$, and $F\ge0$.
\end{proof}

\begin{remark}
The proof shows $\orient(x,\rho(x),\rho^{j}(x))=4\abs{x-c}^2\sin(\pi/m)\sin(j\pi/m)\sin((j-1)\pi/m)$ for
$1\le j\le m$: every vertex of the regular $m$-gon $\{\rho^{j}(x)\}$ lies weakly to the left of the directed edge
$x\to\rho(x)$. (If moreover $x\ne c$ and $m\ge2$, then by (S1) $\sin(\pi/m)>0$, $\sin(j\pi/m)=0$ only for $j=m$ and
$\sin((j-1)\pi/m)=0$ only for $j=1$, so the value is $0$ exactly for $j\in\{1,m\}$; this refinement is not used.)
\end{remark}

\begin{definition}[lin; \leanref{ErdosSzekeres.LayerLemma.lin}{LAYERLEAN/LayerLemma.lean:625}]\label{lg:def:lin}
For points $w,a,y$,
\[
  \operatorname{lin}(w,a,y)=w_0(y_0-a_0)+w_1(y_1-a_1)=\langle w,\,y-a\rangle .
\]
\end{definition}

For $w\ne0$, $\{y:\operatorname{lin}(w,a,y)>0\}$ is the open half-plane bounded by the line through $a$ orthogonal to
$w$, on the side into which $w$ points. Note $\operatorname{lin}(w,a,a)=0$.

\begin{lemma}[the tangent functional on the orbit; \leanref{ErdosSzekeres.LayerLemma.lin_rotAbout}{LAYERLEAN/LayerLemma.lean:794}]\label{lg:lem:lin-rot}
For all points $c,p$ and every $\beta\in\R$,
\[
  \operatorname{lin}\bigl(c-p,\;p,\;\operatorname{rotAbout}(c,\beta)(p)\bigr)=\bigl((p_0-c_0)^2+(p_1-c_1)^2\bigr)(1-\cos\beta),
\]
where $c-p$ is the coordinatewise difference.
\end{lemma}

\begin{proof}
Put $u=p-c$; then $c-p=-u$ and $\operatorname{rotAbout}(c,\beta)(p)-p=R_\beta u-u=(\cos\beta-1)u+\sin\beta\,Ju$. Hence,
by~\eqref{lg:eq:J},
\[
  \operatorname{lin}(c-p,p,\operatorname{rotAbout}(c,\beta)(p))=\langle-u,(\cos\beta-1)u+\sin\beta\,Ju\rangle
  =(1-\cos\beta)\abs{u}^2-\sin\beta\,\langle u,Ju\rangle=(1-\cos\beta)\abs{u}^2.\qedhere
\]
\end{proof}

\begin{lemma}[other vertices lie strictly inside the tangent; \leanref{ErdosSzekeres.LayerLemma.lin_orbit_pos}{LAYERLEAN/LayerLemma.lean:801}]\label{lg:lem:lin-orbit}
Let $c,p$ be points with $p\ne c$, let $m\in\N$ with $0<m$, let $j\in\N$ with $0<j$ and $j<m$, and let
$\rho=\operatorname{rotAbout}(c,2\pi/m)$. Then
\[
  0<\operatorname{lin}\bigl(c-p,\;p,\;\rho^{j}(p)\bigr).
\]
\end{lemma}

\begin{proof}
By Lemma~\ref{lg:lem:rot-iter}, $\rho^{j}(p)=\operatorname{rotAbout}(c,\theta)(p)$ with $\theta=j\cdot(2\pi/m)$, so by
Lemma~\ref{lg:lem:lin-rot} the value is $\abs{p-c}^2(1-\cos\theta)$. Since $p\ne c$, one of $p_0-c_0$, $p_1-c_1$ is
non-zero, so $\abs{p-c}^2>0$. Since $0<j<m$, $0<\theta<2\pi$. By (S2), $\cos\theta\le1$, and $\cos\theta=1$ only for
$\theta\in2\pi\Z$, which has no point in $(0,2\pi)$; hence $1-\cos\theta>0$, and the product is positive.
\end{proof}

\begin{lemma}[next and previous vertex; \leanref{ErdosSzekeres.LayerLemma.orient_next_prev_pos}{LAYERLEAN/LayerLemma.lean:830}]\label{lg:lem:next-prev}
Let $c,p$ be points with $p\ne c$, let $m\in\N$ with $3\le m$, and let $\rho=\operatorname{rotAbout}(c,2\pi/m)$. Then
\[
  0<\orient\bigl(p,\rho(p),\rho^{m-1}(p)\bigr).
\]
\end{lemma}

\begin{proof}
As $m\ge1$, the natural number $m-1$ has real cast $(m:\R)-1$. By Lemma~\ref{lg:lem:rot-iter},
$\rho^{m-1}(p)=\operatorname{rotAbout}(c,\beta)(p)$ with $\beta=(m-1)\cdot(2\pi/m)$. By Lemma~\ref{lg:lem:orient-rot}
with $\alpha=2\pi/m$, the orientation equals $\abs{p-c}^2F$ with $F$ as in the proof of Lemma~\ref{lg:lem:regular}.
Put $t=\pi/m$, so $mt=\pi$, $\alpha=2t$ and $\beta=2((m-1)t)$. Lemma~\ref{lg:lem:trig} with $s=(m-1)t$ gives
\[
  F=4\sin t\,\sin((m-1)t)\,\sin((m-1)t-t)=4\sin t\,\sin(\pi-t)\,\sin(\pi-2t).
\]
Since $m\ge3$, $0<t\le\pi/3$; hence $t$, $\pi-t$ and $\pi-2t$ (which is $\ge\pi/3$) all lie in the open interval
$(0,\pi)$, where the sine is positive by (S1). So $F>0$; also $\abs{p-c}^2>0$ as $p\ne c$. The product is positive.
\end{proof}

\subsection{Convexity of half-planes and disks}\label{lg:ssec:convex}

\begin{lemma}[$\{\orient(a,b,\cdot)\ge0\}$ is convex; \leanref{ErdosSzekeres.LayerLemma.convex_orient_nonneg}{LAYERLEAN/LayerLemma.lean:516}]\label{lg:lem:halfplane-convex}
For all points $a,b$ the set $\{q\in\R^2:0\le\orient(a,b,q)\}$ is convex.
\end{lemma}

\begin{proof}
Let $q_1,q_2$ be in the set and $u,v\ge0$ with $u+v=1$. By (O4),
$\orient(a,b,uq_1+vq_2)=u\orient(a,b,q_1)+v\orient(a,b,q_2)$, a sum of two products of non-negative numbers.
\end{proof}

For $a\ne b$ this set is the closed half-plane to the left of the directed line $a\to b$, the line included; for
$a=b$ it is the whole plane, since $\orient(a,a,q)=0$ by (O3).

\begin{definition}[dist2; \leanref{ErdosSzekeres.LayerLemma.dist2}{LAYERLEAN/LayerLemma.lean:999}]\label{lg:def:dist2}
For points $c,y$, $\operatorname{dist2}(c,y)=(y_0-c_0)^2+(y_1-c_1)^2=\abs{y-c}^2$, the squared Euclidean distance
(defined by this formula, not through a norm).
\end{definition}

\begin{lemma}[disks; \leanref{ErdosSzekeres.LayerLemma.convex_disk}{LAYERLEAN/LayerLemma.lean:1002}]\label{lg:lem:disk}
For every point $c$ and every $R\in\R$ the set $\{y\in\R^2:\operatorname{dist2}(c,y)\le R\}$ is convex.
\end{lemma}

\begin{proof}
Let $y_1,y_2$ be in the set and $u,v\ge0$ with $u+v=1$. Put $A=y_1-c$ and $B=y_2-c$. Since $u+v=1$,
$(uy_1+vy_2)-c=uA+vB$, and, using $1-u=v$ and $1-v=u$,
\[
  u\abs{A}^2+v\abs{B}^2-\abs{uA+vB}^2=u(1-u)\abs{A}^2+v(1-v)\abs{B}^2-2uv\langle A,B\rangle=uv\,\abs{A-B}^2\ge0.
\]
Hence $\operatorname{dist2}(c,uy_1+vy_2)=\abs{uA+vB}^2\le u\abs{A}^2+v\abs{B}^2\le uR+vR=R$. (Lean proves the same
inequality with \lean{nlinarith}, from the non-negativity of the squares of the coordinate differences of $y_1,y_2$ and
of $uv$.)
\end{proof}

For $R\ge0$ the set is the closed disk of radius $\sqrt R$ about $c$; for $R<0$ it is empty.

\subsection{Strict nesting, the centre, and the radius order}\label{lg:ssec:nest}

\begin{theorem}[strict nesting of the convex layers; \leanref{ErdosSzekeres.LayerLemma.strict_nesting}{LAYERLEAN/LayerLemma.lean:529}]\label{lg:thm:nesting}
Let $S$ be a finite set of points, $c$ a point, and $m,K\in\N$ with $2\le m$ and $K\le 2m$; let
$\rho=\operatorname{rotAbout}(c,2\pi/m)$. Assume that $c\notin S$, that $\rho$ maps $S$ into $S$, that $S$ is $K$-free,
and that $S$ is in general position. Let $n,n'\in\N$ with $n<n'$, let $x\in\hv(\peel(S,n))$ and $y\in\peel(S,n')$. Then
\[
  0<\orient\bigl(x,\rho(x),y\bigr).
\]
\end{theorem}

\begin{proof}
Since $m\ge2$, also $0<m$, so Lemma~\ref{lg:lem:layer} applies to $n$ and $x$. Put $L=\hv(\peel(S,n))$. By
Lemma~\ref{lg:lem:layer}(a), $L=\{\rho^{j}(x):0\le j<m\}$; by (c), applied with $n'\ge n$, $y\in\conv(L)$; by (d),
applied with $n'>n$, $y\notin L$.

\emph{Weak inequality.} Let $H=\{q:0\le\orient(x,\rho(x),q)\}$. Every element of $L$ is $\rho^{j}(x)$ with $j<m$, in
particular $j\le m$, so it lies in $H$ by Lemma~\ref{lg:lem:regular}. $H$ is convex by
Lemma~\ref{lg:lem:halfplane-convex} (with $a=x$, $b=\rho(x)$). Hence $\conv(L)\subseteq H$, and $y\in H$, that is,
$0\le\orient(x,\rho(x),y)$.

\emph{Strict inequality.} As $1<m$, $\rho(x)=\rho^{1}(x)\in L$ by Lemma~\ref{lg:lem:layer}(a). The three points
$x,\rho(x),y$ lie in $S$: $x,\rho(x)\in L\subseteq S$ by Lemma~\ref{lg:lem:layer}(b), and $y\in\peel(S,n')\subseteq S$
by (L2). They are pairwise distinct:
\begin{itemize}
\item $x\ne\rho(x)$: $x\ne c$ by Lemma~\ref{lg:lem:layer}(b), so Lemma~\ref{lg:lem:lin-orbit} with $p=x$ and $j=1$
  ($0<1<m$) gives $0<\operatorname{lin}(c-x,x,\rho(x))$, whereas $\operatorname{lin}(c-x,x,x)=0$;
\item $y\ne x$ and $y\ne\rho(x)$, because $x,\rho(x)\in L$ and $y\notin L$.
\end{itemize}
By general position and Lemma~\ref{fd:lem:collinear}, $\orient(x,\rho(x),y)\ne0$. Together with the weak inequality,
$0<\orient(x,\rho(x),y)$.
\end{proof}

\begin{remark}
(1) Lean proves $x\ne\rho(x)$ differently: $x\ne c$ has minimal period $m\ge2$ under $\rho$
(Lemma~\ref{ll:lem:minperiod}), while a fixed point of $\rho$ has minimal period $1$.
(2) For $m=2$ the hypotheses of Theorem~\ref{lg:thm:nesting} are never all satisfied. Indeed, then
$L=\{x,\rho(x)\}$ and $y\in\conv\{x,\rho(x)\}$. The segment $\{(1-\lambda)x+\lambda\rho(x):0\le\lambda\le1\}$ contains
$L$ and is convex (a convex combination of the points with parameters $\lambda_1,\lambda_2$ is the point with the
parameter given by the same combination of $\lambda_1,\lambda_2$), so it contains $\conv\{x,\rho(x)\}$, and
$y=(1-\lambda)x+\lambda\rho(x)$ with $0\le\lambda\le1$. By (O4)
and (O3) $\orient(x,\rho(x),y)=(1-\lambda)\orient(x,\rho(x),x)+\lambda\orient(x,\rho(x),\rho(x))=0$, contradicting the
conclusion. So for $m=2$ the theorem says that, under its hypotheses on $S$, $\peel(S,n')=\emptyset$ for all $n'>n$ as
soon as $\hv(\peel(S,n))\ne\emptyset$.
(3) The theorem is the orientation statement behind the nesting clauses of the layered encoders. That the encoders
emit exactly these clauses is not part of \dirn{LAYERLEAN/LayerLemma.lean}; see the paragraph \emph{Encoders} at the
beginning of this section.
\end{remark}

\begin{lemma}[even orders exclude the centre; \leanref{ErdosSzekeres.LayerLemma.centre_not_mem_of_even}{LAYERLEAN/LayerLemma.lean:582}]\label{lg:lem:centre}
Let $S$ be a finite set, $c$ a point, and $m\in\N$ with $0<m$ and $m$ even; let $\rho=\operatorname{rotAbout}(c,2\pi/m)$.
Assume that $\rho$ maps $S$ into $S$ and that $S$ is in general position. Let $p$ be a point with $p\in S$ and
$p\ne c$. Then $c\notin S$.
\end{lemma}

\begin{proof}
Suppose $c\in S$. Since $m$ is even, $m=k+k$ for some $k\in\N$ (Mathlib's \lean{Even}); $k\ge1$ because $m>0$, and the
natural-number quotient $m/2$ equals $k$. Let $q=\rho^{k}(p)$ (Lean: $\rho^{m/2}(p)$). Since $\rho$ maps $S$ into $S$,
$q\in S$. By Corollary~\ref{lg:cor:period}, $\rho^{k}=\operatorname{rotAbout}(c,2\pi k/m)$, and
$2\pi k/m=2\pi k/(2k)=\pi$. As $\cos\pi=-1$ and $\sin\pi=0$,
\[
  q=\operatorname{rotAbout}(c,\pi)(p)=\bigl(c_0-(p_0-c_0),\;c_1-(p_1-c_1)\bigr)=2c-p .
\]
By (O1), $\orient(p,c,q)=\det(c-p,\,2c-2p)=2\det(c-p,\,c-p)=0$. The points $p,c,q$ are pairwise distinct: $p\ne c$ by
hypothesis; $q=c$ would mean $2c-p=c$, i.e.\ $p=c$; and $p=q$ would mean $p=2c-p$, i.e.\ $p=c$. Since $p,c,q\in S$,
general position and Lemma~\ref{fd:lem:collinear} give $\orient(p,c,q)\ne0$, a contradiction. Hence $c\notin S$.
\end{proof}

\begin{remark}
There is no $K$-freeness hypothesis and no hypothesis $c\notin S$. The hypothesis that $S$ has a point $p\neq c$ cannot be
dropped: $S=\{c\}$ is mapped into itself by $\rho$, is in general position (it has no three distinct points), and
contains $c$.
\end{remark}

\begin{theorem}[the radius order of the layers; \leanref{ErdosSzekeres.LayerLemma.dist2_le_of_later}{LAYERLEAN/LayerLemma.lean:1012}]\label{lg:thm:radius}
Let $S$ be a finite set, $c$ a point, and $m,K\in\N$ with $0<m$ and $K\le 2m$; let $\rho=\operatorname{rotAbout}(c,2\pi/m)$.
Assume that $c\notin S$, that $\rho$ maps $S$ into $S$, and that $S$ is $K$-free. Let $n,n'\in\N$ with $n\le n'$, let
$x\in\hv(\peel(S,n))$ and $y\in\peel(S,n')$. Then
\[
  \operatorname{dist2}(c,y)\le\operatorname{dist2}(c,x).
\]
\end{theorem}

\begin{proof}
Put $L=\hv(\peel(S,n))$. By Lemma~\ref{lg:lem:layer}(a), $L=\{\rho^{j}(x):0\le j<m\}$. For each $j$,
$\rho^{j}(x)=\operatorname{rotAbout}(c,2\pi j/m)(x)$ by Corollary~\ref{lg:cor:period}, so
$\operatorname{dist2}(c,\rho^{j}(x))=\operatorname{dist2}(c,x)$ by Lemma~\ref{lg:lem:rotfacts}(d). Hence $L$ is
contained in $D=\{z:\operatorname{dist2}(c,z)\le\operatorname{dist2}(c,x)\}$, which is convex by Lemma~\ref{lg:lem:disk}
(with $R=\operatorname{dist2}(c,x)$). So $\conv(L)\subseteq D$. By Lemma~\ref{lg:lem:layer}(c) (as $n'\ge n$),
$y\in\conv(L)\subseteq D$.
\end{proof}

\begin{remark}
Only the weak inequality is stated, for $n\le n'$ (including $n=n'$), and general position is not assumed. Strictly
decreasing radii are not claimed here. Rule~B (Lemma~\ref{lg:lem:ruleB}) uses only weak inequalities between radii.
\end{remark}

\subsection{Fans and the chain clauses}\label{lg:ssec:chain}

The two orientation identities behind the fan lemma, and the identity behind Lemma~\ref{lg:lem:halfplane}, are all
instances of one vector identity in the plane.

\begin{lemma}[a Cramer identity; written step]\label{lg:lem:cramer}
For all $X,Y,Z\in\R^2$,
\[
  \det(X,Y)\,Z+\det(Y,Z)\,X+\det(Z,X)\,Y=0 .
\]
Consequently, for all $w,V\in\R^2$,
\begin{align*}
  \det(X,Y)\langle w,Z\rangle+\det(Y,Z)\langle w,X\rangle+\det(Z,X)\langle w,Y\rangle&=0,\\
  \det(X,Y)\det(V,Z)+\det(Y,Z)\det(V,X)+\det(Z,X)\det(V,Y)&=0.
\end{align*}
\end{lemma}

\begin{proof}
Fix $k\in\{0,1\}$ and consider the $3\times3$ matrix with rows $(X_k,Y_k,Z_k)$, $(X_0,Y_0,Z_0)$, $(X_1,Y_1,Z_1)$. Two of
its rows are equal, so its determinant vanishes. Expanding along the first row gives
$X_k\det(Y,Z)-Y_k\det(X,Z)+Z_k\det(X,Y)=0$, which is the $k$-th coordinate of the first identity because
$-\det(X,Z)=\det(Z,X)$. The other two identities follow by applying the linear maps $v\mapsto\langle w,v\rangle$ and
$v\mapsto\det(V,v)$.
\end{proof}

\begin{lemma}[transitivity in a half-plane; \leanref{ErdosSzekeres.LayerLemma.orient_trans_halfplane}{LAYERLEAN/LayerLemma.lean:629}]\label{lg:lem:trans}
Let $a,b,c,e,w$ be points (here $c$ is an arbitrary point, not a centre). If
\[
  0<\operatorname{lin}(w,a,b),\quad 0<\operatorname{lin}(w,a,c),\quad 0<\operatorname{lin}(w,a,e),\quad
  0<\orient(a,b,c),\quad 0<\orient(a,c,e),
\]
then $0<\orient(a,b,e)$.
\end{lemma}

\begin{proof}
Put $B=b-a$, $C=c-a$, $E=e-a$. By (O1) and Definition~\ref{lg:def:lin}, $\orient(a,b,c)=\det(B,C)$,
$\orient(a,c,e)=\det(C,E)$, $\orient(a,b,e)=\det(B,E)$, and $\operatorname{lin}(w,a,z)=\langle w,z-a\rangle$. The
second identity of Lemma~\ref{lg:lem:cramer} with $(X,Y,Z)=(B,C,E)$ reads
$\det(B,C)\langle w,E\rangle+\det(C,E)\langle w,B\rangle+\det(E,B)\langle w,C\rangle=0$. Since $\det(E,B)=-\det(B,E)$,
this is
\begin{equation}\label{lg:eq:trans}
  \orient(a,b,e)\operatorname{lin}(w,a,c)=\orient(a,b,c)\operatorname{lin}(w,a,e)+\orient(a,c,e)\operatorname{lin}(w,a,b).
\end{equation}
The right side is a sum of two products of positive numbers, hence positive. As $\operatorname{lin}(w,a,c)>0$, it
follows that $\orient(a,b,e)>0$. (Lean checks identity~\eqref{lg:eq:trans} directly by \lean{ring}, without
Lemma~\ref{lg:lem:cramer}.)
\end{proof}

In words: inside an open half-plane bounded by a line through $a$, the counterclockwise angular order around $a$ is
transitive.

\begin{lemma}[fan step; \leanref{ErdosSzekeres.LayerLemma.fan_step}{LAYERLEAN/LayerLemma.lean:641}]\label{lg:lem:fan-step}
Let $p,a,b,c,e$ be points (here $c$ is an arbitrary point, not a centre). If
\[
  0<\orient(b,c,e),\quad 0<\orient(p,c,e),\quad 0<\orient(p,b,c),\quad 0<\orient(a,b,c),\quad 0<\orient(p,a,c),
\]
then $0<\orient(a,c,e)$.
\end{lemma}

\begin{proof}
Put $P=p-c$, $A=a-c$, $B=b-c$, $E=e-c$. By (O2) (cyclic invariance, and $\orient(x,y,c)=\det(x-c,y-c)$ for all
$x,y$),
\begin{gather*}
  \orient(p,b,c)=\det(P,B),\qquad \orient(a,b,c)=\det(A,B),\qquad \orient(p,a,c)=\det(P,A),\\
  \orient(a,c,e)=\orient(e,a,c)=\det(E,A),\qquad \orient(p,c,e)=\orient(e,p,c)=\det(E,P),\\
  \orient(b,c,e)=\orient(e,b,c)=\det(E,B).
\end{gather*}
The third identity of Lemma~\ref{lg:lem:cramer} with $(X,Y,Z,V)=(P,B,A,E)$ reads
\[
  \det(P,B)\det(E,A)+\det(B,A)\det(E,P)+\det(A,P)\det(E,B)=0.
\]
Since $\det(B,A)=-\det(A,B)$ and $\det(A,P)=-\det(P,A)$, this is
\begin{equation}\label{lg:eq:fanstep}
  \orient(p,b,c)\orient(a,c,e)=\orient(p,c,e)\orient(a,b,c)+\orient(b,c,e)\orient(p,a,c).
\end{equation}
By the hypotheses the right side is positive, and $\orient(p,b,c)>0$; hence $\orient(a,c,e)>0$. (Lean checks
identity~\eqref{lg:eq:fanstep} directly by \lean{ring}, without Lemma~\ref{lg:lem:cramer}.)
\end{proof}

\begin{lemma}[orientation as a linear functional; \leanref{ErdosSzekeres.LayerLemma.orient_eq_lin}{LAYERLEAN/LayerLemma.lean:653}]\label{lg:lem:orient-lin}
For all points $p,x,y$,
\[
  \orient(p,x,y)=\operatorname{lin}\bigl((-(x_1-p_1),\;x_0-p_0),\;x,\;y\bigr).
\]
\end{lemma}

\begin{proof}
By (O1), $\orient(p,x,y)=\det(x-p,\,y-p)=\det(x-p,\,y-x)+\det(x-p,\,x-p)=\det(x-p,\,y-x)=\langle J(x-p),\,y-x\rangle$,
and $J(x-p)=(-(x_1-p_1),\,x_0-p_0)$.
\end{proof}

So, for $x\ne p$, $\{y:\orient(p,x,y)>0\}$ is the open half-plane to the left of the directed line $p\to x$; its
bounding line passes through $x$.

\begin{lemma}[fan chains have all triples positive; \leanref{ErdosSzekeres.LayerLemma.fan_positive}{LAYERLEAN/LayerLemma.lean:658}]\label{lg:lem:fan-pos}
Let $p,w$ be points, $q:\N\to\R^2$ a sequence of points (write $q_s=q(s)$), and $n\in\N$. Assume
\begin{itemize}
\item[(H)] $0<\operatorname{lin}(w,p,q_s)$ for all $s<n$;
\item[(F)] $0<\orient(p,q_s,q_{s+1})$ for all $s$ with $s+1<n$;
\item[(T)] $0<\orient(q_s,q_{s+1},q_{s+2})$ for all $s$ with $s+2<n$.
\end{itemize}
Then (a) $0<\orient(p,q_j,q_k)$ for all $j,k$ with $j<k$ and $k<n$; and (b) $0<\orient(q_i,q_j,q_k)$ for all $i,j,k$
with $i<j$, $j<k$ and $k<n$.
\end{lemma}

\begin{proof}
(a) Fix $j$. We show by induction on $k\ge j+1$ that $k<n$ implies $\orient(p,q_j,q_k)>0$. For $k=j+1$ this is (F)
with $s=j$. Let $k\ge j+1$ satisfy the claim and let $k+1<n$. Then $k<n$, so $\orient(p,q_j,q_k)>0$; (F) with $s=k$
gives $\orient(p,q_k,q_{k+1})>0$; and (H) gives $\operatorname{lin}(w,p,q_t)>0$ for $t\in\{j,k,k+1\}$, all of which
are $<n$. Lemma~\ref{lg:lem:trans} with $(a,b,c,e)=(p,q_j,q_k,q_{k+1})$ gives $\orient(p,q_j,q_{k+1})>0$.

(b$_0$) We first show: for all $i<s$ with $s+1<n$, $\orient(q_i,q_s,q_{s+1})>0$. Fix $i$ and induct on $s\ge i+1$. For
$s=i+1$ this is (T) with $s=i$ (note $i+2<n$). Let $s\ge i+1$ satisfy the claim and let $s+2<n$. Apply
Lemma~\ref{lg:lem:fan-step} with $(p,a,b,c,e)=(p,q_i,q_s,q_{s+1},q_{s+2})$; its five hypotheses hold:
$\orient(q_s,q_{s+1},q_{s+2})>0$ by (T); $\orient(p,q_{s+1},q_{s+2})>0$ by (F) with $s+1$; $\orient(p,q_s,q_{s+1})>0$ by
(F) with $s$; $\orient(q_i,q_s,q_{s+1})>0$ by the induction hypothesis (as $s+1<n$); and $\orient(p,q_i,q_{s+1})>0$ by
(a), since $i<s+1<n$. The conclusion is $\orient(q_i,q_{s+1},q_{s+2})>0$.

(b) Fix $i<j$ and induct on $k\ge j+1$, showing that $k<n$ implies $\orient(q_i,q_j,q_k)>0$. For $k=j+1$ this is
(b$_0$) with $s=j$. Let $k\ge j+1$ satisfy the claim and let $k+1<n$. Put $w'=(-(q_{i,1}-p_1),\,q_{i,0}-p_0)$, where
$q_i=(q_{i,0},q_{i,1})$. By Lemma~\ref{lg:lem:orient-lin}, $\operatorname{lin}(w',q_i,y)=\orient(p,q_i,y)$ for every $y$,
so by (a) $\operatorname{lin}(w',q_i,q_t)>0$ for every $t$ with $i<t<n$, in particular for $t\in\{j,k,k+1\}$. Moreover
$\orient(q_i,q_j,q_k)>0$ by the induction hypothesis (as $k<n$), and $\orient(q_i,q_k,q_{k+1})>0$ by (b$_0$) with
$s=k$ (as $i<k$ and $k+1<n$). Lemma~\ref{lg:lem:trans} with $(a,b,c,e,w)=(q_i,q_j,q_k,q_{k+1},w')$ gives
$\orient(q_i,q_j,q_{k+1})>0$.
\end{proof}

\begin{definition}[fan points; \leanref{ErdosSzekeres.LayerLemma.fanPt}{LAYERLEAN/LayerLemma.lean:694}]\label{lg:def:fanpt}
For $n\in\N$, a point $p$, $q:\N\to\R^2$ and $i\in\{0,1,\dots,n\}$ (Lean \texttt{Fin (n + 1)}; $n$ is an implicit
argument),
\[
  \operatorname{fanPt}(p,q)(i)=\begin{cases} p & \text{if } i=0,\\ q(i-1) & \text{otherwise.}\end{cases}
\]
\end{definition}

\begin{theorem}[fan chains give convex polygons; \leanref{ErdosSzekeres.LayerLemma.hasConvexNGon_of_fan}{LAYERLEAN/LayerLemma.lean:700}]\label{lg:thm:fan}
Let $S$ be a finite set, $p,w$ points, $q:\N\to\R^2$ (write $q_s=q(s)$) and $n\in\N$ with $3\le n$. Assume that $p\in S$,
that $q_s\in S$ for all $s<n$, and that (H), (F) and (T) of Lemma~\ref{lg:lem:fan-pos} hold:
$0<\operatorname{lin}(w,p,q_s)$ for all $s<n$; $0<\orient(p,q_s,q_{s+1})$ for all $s$ with $s+1<n$; and
$0<\orient(q_s,q_{s+1},q_{s+2})$ for all $s$ with $s+2<n$. Then \texttt{HasConvexNGon (n + 1) S}: some $(n+1)$-element
subset of $S$ is in convex position.
\end{theorem}

\begin{proof}
Let $x_i=\operatorname{fanPt}(p,q)(i)$ for $i\in\{0,\dots,n\}$; so $x_0=p$ and $x_i=q_{i-1}$ for $1\le i\le n$.

\emph{Step 1: every increasing triple is positive.} Let $0\le i<j<k\le n$. If $i=0$, then $j,k\ge1$ and
$\orient(x_0,x_j,x_k)=\orient(p,q_{j-1},q_{k-1})>0$ by Lemma~\ref{lg:lem:fan-pos}(a), since $j-1<k-1<n$. If $i\ge1$,
then $\orient(x_i,x_j,x_k)=\orient(q_{i-1},q_{j-1},q_{k-1})>0$ by Lemma~\ref{lg:lem:fan-pos}(b), since
$i-1<j-1<k-1<n$.

\emph{Step 2: the points $x_0,\dots,x_n$ are pairwise distinct.} Suppose $x_i=x_j$ with $i\ne j$. As $n+1\ge4>2$, there
is $k\in\{0,\dots,n\}$ different from $i$ and $j$. Arrange $i,j,k$ increasingly as $i'<j'<k'$. By Step~1,
$\orient(x_{i'},x_{j'},x_{k'})>0$; but two of these three points are equal (namely $x_i$ and $x_j$), so the orientation
is $0$ by (O3), a contradiction. (Lean chooses $k$ explicitly and treats the possible orders of $i,j,k$ one by one,
using \lean{orient_cycle} and $\orient(a,a,d)=0$; the argument is the same.)

\emph{Step 3: convex position.} Apply Lemma~\ref{fd:lem:pos-triples} (\lean{Suk.convexPosition_of_positive_triples})
with the linearly ordered index set $\{0,\dots,n\}$ (Lean \texttt{Fin (n + 1)} with its usual order), the injective map
$i\mapsto x_i$ (Step~2), and the full index set, which has $n+1\ge4$ elements; its hypothesis is Step~1. Hence
$T=\{x_0,\dots,x_n\}$ is in convex position.

\emph{Step 4.} $T\subseteq S$, since $x_0=p\in S$ and $x_i=q_{i-1}\in S$ for $1\le i\le n$ (as $i-1<n$); and
$\card T=n+1$ by Step~2. Hence \texttt{HasConvexNGon (n + 1) S}.
\end{proof}

\begin{remark}
General position is not assumed, and the chain points are not assumed distinct: distinctness of $p,q_0,\dots,q_{n-1}$
follows from the positivity of the orientations (Step~2), and the proof shows that $\{p,q_0,\dots,q_{n-1}\}$ itself is
the convex $(n+1)$-gon.
\end{remark}

\begin{lemma}[inner side of the tangent in the layers; \leanref{ErdosSzekeres.LayerLemma.lin_pos_of_peel}{LAYERLEAN/LayerLemma.lean:864}]\label{lg:lem:halfplane}
Let $S$ be a finite set, $c$ a point, and $m,K\in\N$ with $3\le m$ and $K\le 2m$; let $\rho=\operatorname{rotAbout}(c,2\pi/m)$.
Assume that $c\notin S$, that $\rho$ maps $S$ into $S$, that $S$ is $K$-free, and that $S$ is in general position. Let
$i\in\N$, $p\in\hv(\peel(S,i))$, and $y\in\peel(S,i)$ with $y\ne p$. Then
\[
  0<\operatorname{lin}(c-p,\;p,\;y).
\]
\end{lemma}

\begin{proof}
Since $m\ge3$, $0<m$ and Lemma~\ref{lg:lem:layer} applies to $i$ and $p$. Put $L=\hv(\peel(S,i))$; then
$L=\{\rho^{j}(p):0\le j<m\}$ and $p\ne c$.

\emph{Case $y\in L$.} Then $y=\rho^{j}(p)$ for some $j<m$, and $j\ne0$ because $\rho^{0}(p)=p\ne y$. So $0<j<m$, and
Lemma~\ref{lg:lem:lin-orbit} gives $0<\operatorname{lin}(c-p,p,y)$.

\emph{Case $y\notin L$.} Then $y\in\peel(S,i)\setminus L=\peel(S,i+1)$. Let $p'=\rho^{m-1}(p)$. Since $m-1<m$,
$p'\in L$ by Lemma~\ref{lg:lem:layer}(a); and $\rho(p')=\rho^{m}(p)=p$ by Corollary~\ref{lg:cor:period}. We collect
five positive quantities.
\begin{itemize}
\item[(e1)] $\orient(p,\rho(p),y)>0$: Theorem~\ref{lg:thm:nesting} with $(n,n',x)=(i,i+1,p)$ (its hypotheses hold:
  $2\le m$, $K\le2m$, $c\notin S$, invariance, $K$-freeness, general position, $p\in\hv(\peel(S,i))$,
  $y\in\peel(S,i+1)$).
\item[(e2)] $\orient(p',p,y)>0$: Theorem~\ref{lg:thm:nesting} with $(n,n',x)=(i,i+1,p')$ gives
  $\orient(p',\rho(p'),y)>0$, and $\rho(p')=p$.
\item[(w1)] $\operatorname{lin}(c-p,p,\rho(p))>0$: Lemma~\ref{lg:lem:lin-orbit} with $j=1$ (as $0<1<m$).
\item[(w2)] $\operatorname{lin}(c-p,p,p')>0$: Lemma~\ref{lg:lem:lin-orbit} with $j=m-1$ (as $0<m-1<m$).
\item[(o)] $\orient(p,\rho(p),p')>0$: Lemma~\ref{lg:lem:next-prev} (as $p\ne c$ and $m\ge3$).
\end{itemize}
Now put $B=\rho(p)-p$, $C=p'-p$, $E=y-p$ and $w=c-p$. By (O1) and (O2),
$\orient(p,\rho(p),p')=\det(B,C)$, $\orient(p,\rho(p),y)=\det(B,E)$ and $\orient(p',p,y)=\orient(p,y,p')=\det(E,C)$;
and $\operatorname{lin}(w,p,z)=\langle w,z-p\rangle$. The second identity of Lemma~\ref{lg:lem:cramer} with
$(X,Y,Z)=(B,C,E)$ gives $\det(B,C)\langle w,E\rangle=-\det(C,E)\langle w,B\rangle-\det(E,B)\langle w,C\rangle
=\det(E,C)\langle w,B\rangle+\det(B,E)\langle w,C\rangle$, that is,
\begin{equation}\label{lg:eq:key}
  \orient(p,\rho(p),p')\operatorname{lin}(c-p,p,y)
  =\orient(p',p,y)\operatorname{lin}(c-p,p,\rho(p))+\orient(p,\rho(p),y)\operatorname{lin}(c-p,p,p').
\end{equation}
By (e2), (w1), (e1) and (w2) the right side is positive, and by (o) $\orient(p,\rho(p),p')>0$. Hence
$\operatorname{lin}(c-p,p,y)>0$.
\end{proof}

\begin{remark}
(1) Geometrically, $\{y:\operatorname{lin}(c-p,p,y)=0\}$ is the tangent at $p$ to the circle about $c$ through $p$, and
the lemma says that every other point of the layers $\ge i$ lies strictly on the side of the centre. Identity
\eqref{lg:eq:key} holds for arbitrary points $p,\rho(p),p',y,c$; Lean proves it by \lean{ring}.
(2) Lean obtains $\rho(p')=\rho^{m}(p)=p$ from the minimal period of $p$ (Lemma~\ref{ll:lem:minperiod}); we used
Corollary~\ref{lg:cor:period} instead.
\end{remark}

\begin{theorem}[geometric chain soundness; \leanref{ErdosSzekeres.LayerLemma.chain_clauses_sound}{LAYERLEAN/LayerLemma.lean:914}]\label{lg:thm:chain}
Let $S$ be a finite set, $c$ a point, and $m,K\in\N$ with $3\le m$, $4\le K$ and $K\le 2m$; let
$\rho=\operatorname{rotAbout}(c,2\pi/m)$. Assume that $c\notin S$, that $\rho$ maps $S$ into $S$, that $S$ is $K$-free,
and that $S$ is in general position. Let $i\in\N$, $p\in\hv(\peel(S,i))$, and $q:\N\to\R^2$ (write $q_s=q(s)$) be such
that
\begin{itemize}
\item[(C1)] $q_s\in\peel(S,i)$ for all $s<K-1$;
\item[(C2)] $q_s\ne p$ for all $s<K-1$;
\item[(C3)] $0<\orient(p,q_s,q_{s+1})$ for all $s$ with $s+1<K-1$;
\item[(C4)] $0<\orient(q_s,q_{s+1},q_{s+2})$ for all $s$ with $s+2<K-1$.
\end{itemize}
Then a contradiction follows (the Lean conclusion is \lean{False}). In other words, no such $i$, $p$ and $q$ exist.
\end{theorem}

\begin{proof}
Since $K\ge4$, $K-1\ge3$ is the ordinary difference. We have $p\in\peel(S,i)$, as $\hv(\peel(S,i))\subseteq\peel(S,i)$.
For every $s<K-1$, Lemma~\ref{lg:lem:halfplane} with $y=q_s$ (which lies in $\peel(S,i)$ by (C1) and differs from $p$
by (C2)) gives $0<\operatorname{lin}(c-p,p,q_s)$. Apply Theorem~\ref{lg:thm:fan} to the finite set $\peel(S,i)$, the
pivot $p$, the vector $w=c-p$, the sequence $q$ and $n=K-1\ge3$: its hypotheses are $p\in\peel(S,i)$, (C1), the
half-plane inequalities just proved, (C3) and (C4). Hence $\peel(S,i)$ contains $(K-1)+1=K$ points in convex position.
Since $\peel(S,i)\subseteq S$ by (L2), so does $S$ (Lemma~\ref{fd:lem:mono}(c)), contradicting the $K$-freeness of
$S$. (Lean unfolds the definition of \lean{HasConvexNGon} here instead of citing Lemma~\ref{fd:lem:mono}(c).)
\end{proof}

\begin{remark}
The chain points $q_s$ are not assumed distinct, and they may lie in any layer $\ge i$. The theorem is the geometric
statement behind the chain clauses of the layered encoders (``no pivot of layer $i$ with $K-1$ chain points whose
consecutive fan and turn orientations are positive''). That the encoders' clauses express exactly this family
(including, where auxiliary chain variables are used, that their least-fixpoint values satisfy every chain clause;
\dirn{LEANEXT/review/REPORT.md}, finding F5) is a code-level step outside \dirn{LAYERLEAN/LayerLemma.lean}; see the
paragraph \emph{Encoders} at the beginning of this section.
\end{remark}

\subsection{Polar coordinates and the angular rules}\label{lg:ssec:polar}

\begin{definition}[polar; \leanref{ErdosSzekeres.LayerLemma.polar}{LAYERLEAN/LayerLemma.lean:932}]\label{lg:def:polar}
For a point $c$ and $r,\theta\in\R$,
\[
  \operatorname{polar}(c,r,\theta)=\bigl(c_0+r\cos\theta,\;c_1+r\sin\theta\bigr).
\]
There is no sign condition on $r$.
\end{definition}

\begin{lemma}[orientation of three points in polar form; \leanref{ErdosSzekeres.LayerLemma.orient_polar}{LAYERLEAN/LayerLemma.lean:935}]\label{lg:lem:polar}
For every point $c$ and all $r_a,r_b,r_d,\alpha,\beta,\delta\in\R$,
\[
  \orient\bigl(\operatorname{polar}(c,r_a,\alpha),\operatorname{polar}(c,r_b,\beta),\operatorname{polar}(c,r_d,\delta)\bigr)
  =r_ar_b\sin(\beta-\alpha)+r_br_d\sin(\delta-\beta)+r_dr_a\sin(\alpha-\delta).
\]
\end{lemma}

\begin{proof}
Write $e(\varphi)=(\cos\varphi,\sin\varphi)$ and let $a,b,d$ be the three points; then $a-c=A:=r_ae(\alpha)$,
$b-c=B:=r_be(\beta)$, $d-c=D:=r_de(\delta)$, so $b-a=B-A$ and $d-a=D-A$. By (O1) and bilinearity and antisymmetry of
$\det$,
\begin{align*}
  \orient(a,b,d)=\det(B-A,D-A)&=\det(B,D)-\det(B,A)-\det(A,D)+\det(A,A)\\
  &=\det(A,B)+\det(B,D)+\det(D,A).
\end{align*}
Finally, for all $r,s,\varphi,\psi\in\R$,
$\det(re(\varphi),se(\psi))=rs(\cos\varphi\sin\psi-\sin\varphi\cos\psi)=rs\sin(\psi-\varphi)$; applied to the three
terms this gives the formula. (Lean expands the three sines on the right by the subtraction formula and concludes with
\lean{ring}.)
\end{proof}

\begin{lemma}[rule A, centre inside; \leanref{ErdosSzekeres.LayerLemma.orient_polar_centre_inside}{LAYERLEAN/LayerLemma.lean:943}]\label{lg:lem:ruleA}
(Angular rule A: centre inside the triangle.)
Let $c$ be a point and $r_a,r_b,r_d,\alpha,\beta,\delta\in\R$ with $0<r_a$, $0<r_b$, $0<r_d$, and
\[
  0<\beta-\alpha<\pi,\qquad 0<\delta-\beta<\pi,\qquad 0<\alpha-\delta+2\pi<\pi .
\]
Then $0<\orient\bigl(\operatorname{polar}(c,r_a,\alpha),\operatorname{polar}(c,r_b,\beta),\operatorname{polar}(c,r_d,\delta)\bigr)$.
\end{lemma}

\begin{proof}
By Lemma~\ref{lg:lem:polar} the orientation is $r_ar_b\sin(\beta-\alpha)+r_br_d\sin(\delta-\beta)+r_dr_a\sin(\alpha-\delta)$.
By (S1) the sine is positive on $(0,\pi)$, so $\sin(\beta-\alpha)>0$, $\sin(\delta-\beta)>0$, and, by
$2\pi$-periodicity, $\sin(\alpha-\delta)=\sin(\alpha-\delta+2\pi)>0$. With $r_a,r_b,r_d>0$ all three terms are
positive.
\end{proof}

The hypotheses say $\alpha<\beta<\delta<\alpha+2\pi$ with all three counterclockwise gaps $\beta-\alpha$,
$\delta-\beta$, $\alpha+2\pi-\delta$ (which sum to $2\pi$) in $(0,\pi)$. The name of the rule refers to the position of
$c$ inside the triangle; that geometric fact is neither stated nor used here, and only the orientation statement is
proved. Lean states only the positive orientation.

\begin{lemma}[sine of a sum below the sum of the sines; \leanref{ErdosSzekeres.LayerLemma.sin_add_lt}{LAYERLEAN/LayerLemma.lean:956}]\label{lg:lem:sinadd}
For all $x,y\in\R$ with $0<x<\pi$ and $0<y<\pi$: $\sin(x+y)<\sin x+\sin y$.
\end{lemma}

\begin{proof}
By (S1), $\sin x>0$ and $\sin y>0$. Moreover $\cos x<1$: by (S2), $\cos x\le1$, and $\cos x=1$ only for
$x\in2\pi\Z$, which has no point in $(0,\pi)$; likewise $\cos y<1$. Hence $\sin x\cos y<\sin x$ and $\cos x\sin y<\sin y$, and
$\sin(x+y)=\sin x\cos y+\cos x\sin y<\sin x+\sin y$.
\end{proof}

\begin{lemma}[rule B, middle point outermost; \leanref{ErdosSzekeres.LayerLemma.orient_polar_halfplane}{LAYERLEAN/LayerLemma.lean:974}]\label{lg:lem:ruleB}
(Angular rule B: open half-plane, angularly middle point outermost.)
Let $c$ be a point and $r_a,r_b,r_d,\alpha,\beta,\delta\in\R$ with
\[
  0<r_a,\quad 0<r_d,\quad r_a\le r_b,\quad r_d\le r_b,\quad 0<\beta-\alpha,\quad 0<\delta-\beta,\quad \delta-\alpha<\pi .
\]
Then $0<\orient\bigl(\operatorname{polar}(c,r_a,\alpha),\operatorname{polar}(c,r_b,\beta),\operatorname{polar}(c,r_d,\delta)\bigr)$.
\end{lemma}

\begin{proof}
Let $a,b,d$ denote the three points $\operatorname{polar}(c,r_a,\alpha)$, $\operatorname{polar}(c,r_b,\beta)$,
$\operatorname{polar}(c,r_d,\delta)$. Put $x=\beta-\alpha>0$ and $y=\delta-\beta>0$. Then $x+y=\delta-\alpha<\pi$, so
$x=(x+y)-y<\pi$ and $y=(x+y)-x<\pi$. Since $\alpha-\delta=-(x+y)$ and the sine is odd, Lemma~\ref{lg:lem:polar} gives
\begin{align*}
  \orient(a,b,d)&=r_ar_b\sin x+r_br_d\sin y-r_dr_a\sin(x+y)\\
  &=r_a(r_b-r_d)\sin x+r_d(r_b-r_a)\sin y+r_ar_d\bigl(\sin x+\sin y-\sin(x+y)\bigr).
\end{align*}
Here $\sin x,\sin y>0$ by (S1) (arguments in $(0,\pi)$), $r_a,r_d>0$ and $r_b-r_d,\,r_b-r_a\ge0$, so the first two terms are
$\ge0$; the third is $>0$ by Lemma~\ref{lg:lem:sinadd} and $r_ar_d>0$. Hence the orientation is positive. (Lean uses
$r_ar_b\sin x\ge r_ar_d\sin x$, $r_br_d\sin y\ge r_ar_d\sin y$ and Lemma~\ref{lg:lem:sinadd}, and concludes with
\lean{nlinarith}; this is the same estimate.)
\end{proof}

\begin{remark}
The radius hypotheses are weak inequalities ($r_b>0$ follows from $0<r_a\le r_b$). This is the form of the radius
order of Theorem~\ref{lg:thm:radius}, which gives $\operatorname{dist2}(c,y)\le\operatorname{dist2}(c,x)$, that is,
$r_y\le r_x$ for the radii $r=\sqrt{\operatorname{dist2}(c,\cdot)}$. The rules apply to points given in polar form
with explicit angle representatives. No Lean statement of this file writes the points of a layered configuration in
that form (a point $y$ as $\operatorname{polar}(c,r,\theta)$ with $r^2=\operatorname{dist2}(c,y)$, and a layer as
$\{\operatorname{polar}(c,r_i,\theta_i+2\pi k/m)\}$; \dirn{LEANEXT/review/REPORT.md}, finding F2); that step, and the
choice of which rule applies to which triple, are outside these statements (see the paragraph \emph{Encoders} at the
beginning of this section).
\end{remark}

\paragraph{Coverage.}
All 27 declarations of the scope of this section, which are all the declarations of lines 465--1032 of
\dirn{LAYERLEAN/LayerLemma.lean} (23 theorems and 4 definitions), are covered:
\lean{rotAbout_iterate} (Lemma~\ref{lg:lem:rot-iter}; the section on the layer lemma restates it as well),
\lean{orient_rotAbout} (Lemma~\ref{lg:lem:orient-rot}),
\lean{trig_layer} (Lemma~\ref{lg:lem:trig}, stated just before Lemma~\ref{lg:lem:regular}, whose proof uses it),
\lean{orient_regular_nonneg} (Lemma~\ref{lg:lem:regular}),
\lean{lin} (Definition~\ref{lg:def:lin}),
\lean{lin_rotAbout} (Lemma~\ref{lg:lem:lin-rot}),
\lean{lin_orbit_pos} (Lemma~\ref{lg:lem:lin-orbit}),
\lean{orient_next_prev_pos} (Lemma~\ref{lg:lem:next-prev}),
\lean{convex_orient_nonneg} (Lemma~\ref{lg:lem:halfplane-convex}),
\lean{dist2} (Definition~\ref{lg:def:dist2}),
\lean{convex_disk} (Lemma~\ref{lg:lem:disk}),
\lean{strict_nesting} (Theorem~\ref{lg:thm:nesting}),
\lean{centre_not_mem_of_even} (Lemma~\ref{lg:lem:centre}),
\lean{dist2_le_of_later} (Theorem~\ref{lg:thm:radius}),
\lean{orient_trans_halfplane} (Lemma~\ref{lg:lem:trans}),
\lean{fan_step} (Lemma~\ref{lg:lem:fan-step}),
\lean{orient_eq_lin} (Lemma~\ref{lg:lem:orient-lin}),
\lean{fan_positive} (Lemma~\ref{lg:lem:fan-pos}),
\lean{fanPt} (Definition~\ref{lg:def:fanpt}),
\lean{hasConvexNGon_of_fan} (Theorem~\ref{lg:thm:fan}),
\lean{lin_pos_of_peel} (Lemma~\ref{lg:lem:halfplane}),
\lean{chain_clauses_sound} (Theorem~\ref{lg:thm:chain}),
\lean{polar} (Definition~\ref{lg:def:polar}),
\lean{orient_polar} (Lemma~\ref{lg:lem:polar}),
\lean{orient_polar_centre_inside} (Lemma~\ref{lg:lem:ruleA}),
\lean{sin_add_lt} (Lemma~\ref{lg:lem:sinadd}) and
\lean{orient_polar_halfplane} (Lemma~\ref{lg:lem:ruleB}); all in \dirn{LAYERLEAN/LayerLemma.lean}, namespace
\lean{ErdosSzekeres.LayerLemma}.
The written steps (Lemmas~\ref{lg:lem:layer}, \ref{lg:lem:rotfacts}, \ref{lg:lem:cramer} and
Corollary~\ref{lg:cor:period}) are not Lean declarations; the paragraph \emph{Rigor} at the beginning of the section
says how each relates to the Lean proofs.
Deliberately not covered here:
\begin{itemize}
\item \lean{peel_antitone} (line 460) and \lean{card_peel} (line 1037, docstring at line 1036), which adjoin these
  lines but belong to the section on the layer lemma (Lemma~\ref{ll:lem:antitone}, Theorem~\ref{ll:thm:card-peel});
\item the verbatim copies of these declarations in other units, which are duplicates. The following eight files
  contain lines 11--1147 of \dirn{LAYERLEAN/LayerLemma.lean} verbatim (checked with \texttt{diff}); a search of the
  uncompressed files of the repository finds no further copy:\newline
  \dirn{LEANEXT/capstone/Capstone.lean};\newline
  \dirn{LEANEXT/capstone/Capstone2.lean};\newline
  \dirn{LEANEXT/merged/Merged.lean};\newline
  \dirn{LEANEXT/linked/Linked.lean};\newline
  \dirn{LEANEXT/es7corr/ES7Corollaries.lean};\newline
  \dirn{LEANEXT/review/VacuityTest.lean} (all 1148 lines, followed by the tests of the next item);\newline
  \dirn{ENCBLEAN/geo/ES8Geo/Capstone2.lean};\newline
  \dirn{LAYERLEAN/logs/axioms_recheck_20261106/LayerLemmaAxioms.lean};
\item the 19 declarations of \dirn{LEANEXT/review/tests_tail.lean} (namespace
  \lean{ErdosSzekeres.LayerLemma.Review}), which no section covers: the definitions \lean{L4}, \lean{S4}, \lean{H6} and
  the theorems \lean{mem_S4}, \lean{card_S4}, \lean{angle4}, \lean{rot4}, \lean{hc4}, \lean{inv4}, \lean{gp4},
  \lean{inner4}, \lean{hno4}, \lean{hv4}, \lean{peel1_4}, \lean{nest4}, \lean{orbit4}, \lean{layers4}, \lean{chain4}
  and \lean{inv6}. They are non-vacuity tests, not results (\dirn{LEANEXT/review/REPORT.md}, item~6). For the eight
  points $S_4=\{(\pm2,0),(0,\pm2),(\pm1/2,\pm1/2)\}$, the centre $c=0$, $m=4$ and $K=8$, the theorems \lean{hc4},
  \lean{inv4}, \lean{gp4} and \lean{hno4} state that $c\notin S_4$, that $\operatorname{rotAbout}(0,2\pi/4)$ maps $S_4$
  into $S_4$, that $S_4$ is in general position and that $S_4$ is $8$-free; \lean{hv4} and \lean{peel1_4} state that
  $(2,0)\in\hv(S_4)$ and $(1/2,1/2)\in\peel(S_4,1)$; and \lean{nest4}, \lean{orbit4}, \lean{layers4} and \lean{chain4}
  apply \lean{strict_nesting} (Theorem~\ref{lg:thm:nesting}), \lean{hullVerts_eq_orbitFS_rotAbout}
  (Theorem~\ref{ll:thm:core-rot}), \lean{layers_are_orbits} (Theorem~\ref{ll:thm:layers}) and
  \lean{chain_clauses_sound} (Theorem~\ref{lg:thm:chain}) to this instance. \lean{inv6} states that
  $\operatorname{rotAbout}(0,2\pi/6)$ maps the set \lean{H6} of the six points $(\cos(k\pi/3),\sin(k\pi/3))$,
  $0\le k<6$ (a regular hexagon), into itself. \dirn{LEANEXT/review/VacuityTest.lean} consists of all 1148 lines of
  \dirn{LAYERLEAN/LayerLemma.lean} followed by the 149 lines of \dirn{LEANEXT/review/tests_tail.lean}, as lines
  1149--1297 (both parts checked with \texttt{cmp}), so the same 19 declarations stand there at lines 1160--1276. Both
  files end with \texttt{\#print axioms}, \texttt{\#check} and \texttt{\#print} commands, which are not declarations;
\item declarations in other namespaces that share a short name with one of the 27 (in
  \lean{ErdosSzekeres.OrbitConsistency}, \lean{PhaseClassification}, \lean{C7} and
  \lean{ErdosSzekeres.SymReduction.Review}); they are separate Lean declarations. Those of
  \dirn{LEANEXT/phase/PhaseClassification.lean} (\lean{polar}, \lean{orient_polar}, \lean{orient_polar_centre_inside},
  \lean{sin_add_lt} and \lean{orient_polar_halfplane}) are treated in the section on phase patterns.
\end{itemize}
